% GENERATED FILE. Edit canonical lessons, then run npm run generate:books.
% canonical-source-hash: fnv1a64:7f551d396de5b83d
% canonical-lessons: KA-C528-sneha, KA-C528-bandhavya, KA-C528-odanata, KA-C528-atmiya, KA-C528-raji-madiko

\chapter{Friendship}
\label{ch:ka-friendship}
\hlchaptermodality{528}

\begin{chapteropening}
\textbf{By the end of this chapter:} \emph{I can say friendship, a bond, companionship, dear and to make up.}

Complete the last lesson of chapter 528: I can say friendship, a bond, companionship, dear and to make up.
\end{chapteropening}

\section[\texorpdfstring{\kn{ಸ್ನೇಹ} (snēha)}{ಸ್ನೇಹ (snēha)}]{\kn{ಸ್ನೇಹ} (snēha) — friendship}

\label{lesson:KA-C528-sneha}



\begin{quote}
Before the new one: say the Kannada for crawling on all fours, then the Kannada for a stepmother.
\end{quote}

\subsection*{You'll want to know: \kn{ಸ್ನೇಹ}}
\textbf{\kn{ಸ್ನೇಹ}} — \emph{snēha} — \textquotedblleft{}friendship\textquotedblright{}.

\begin{quote}
\kn{ನಮ್ಮ} \kn{ಸ್ನೇಹ} \kn{ಶಾಲೆಯಲ್ಲಿ} \kn{ಶುರುವಾಯಿತು}.
\end{quote}

\begin{quote}
\emph{namma snēha śāleyalli śuruvāyitu.}
\end{quote}

\begin{quote}
Our friendship began at school.
\end{quote}

\kn{ಸ್ನೇಹಿತ}, a friend, is built on it.

\begin{cousinweb}
\textbf{\kn{ಸ್ನೇಹ}} (\emph{snēha}) is borrowed from Sanskrit \emph{sneha}, \textquotedblleft{}oiliness, smoothness\textquotedblright{}, and from there \textquotedblleft{}tenderness, friendship\textquotedblright{} (Monier-Williams).
\end{cousinweb}

\begin{grammarlens}[title={What you've built}]
One more word for describing what you lived and felt.
\end{grammarlens}

\subsection*{Your turn}
\begin{itemize}
\raggedright
  \item \emph{Say it:} \emph{snēha}
  \item \emph{Say it:} \emph{snēha}, once more
  \item \emph{Say it:} \emph{namma snēha śāleyalli śuruvāyitu.} — Our friendship began at school.
  \item \emph{Recall:} say the Kannada for crawling on all fours, then the Kannada for a stepmother, then say \emph{snēha} again
\end{itemize}

\begin{tcolorbox}[breakable,colback=teal!4,colframe=teal!35!black,title={Before you move on}]
What does \kn{ಸ್ನೇಹ} mean? (\textquotedblleft{}Friendship\textquotedblright{}.) Say it once more.
\end{tcolorbox}
\section[\texorpdfstring{\kn{ಬಾಂಧವ್ಯ} (bāndhavya)}{ಬಾಂಧವ್ಯ (bāndhavya)}]{\kn{ಬಾಂಧವ್ಯ} (bāndhavya) — a bond; close ties}

\label{lesson:KA-C528-bandhavya}



\begin{quote}
Before the new one: say the Kannada for a stepmother, then the Kannada for friendship.
\end{quote}

\subsection*{You'll want to know: \kn{ಬಾಂಧವ್ಯ}}
\textbf{\kn{ಬಾಂಧವ್ಯ}} — \emph{bāndhavya} — \textquotedblleft{}a bond; close ties\textquotedblright{}.

\begin{quote}
\kn{ಎರಡು} \kn{ಕುಟುಂಬಗಳ} \kn{ನಡುವೆ} \kn{ಒಳ್ಳೆಯ} \kn{ಬಾಂಧವ್ಯ} \kn{ಇದೆ}.
\end{quote}

\begin{quote}
\emph{eraḍu kuṭumbagaḷa naḍuve oḷḷeya bāndhavya ide.}
\end{quote}

\begin{quote}
There is a good bond between the two families.
\end{quote}

Related to \kn{ಬಂಧು}, kin.

\begin{cousinweb}
\textbf{\kn{ಬಾಂಧವ್ಯ}} (\emph{bāndhavya}) is a learned abstract noun built on Sanskrit \emph{bāndhava}, \textquotedblleft{}kinsman\textquotedblright{} (CDIAL 9208): the state of being bound as kin.
\end{cousinweb}

\begin{grammarlens}[title={What you've built}]
One more word for describing what you lived and felt.
\end{grammarlens}

\subsection*{Your turn}
\begin{itemize}
\raggedright
  \item \emph{Say it:} \emph{bāndhavya}
  \item \emph{Say it:} \emph{bāndhavya}, once more
  \item \emph{Say it:} \emph{eraḍu kuṭumbagaḷa naḍuve oḷḷeya bāndhavya ide.} — There is a good bond between the two families.
  \item \emph{Recall:} say the Kannada for a stepmother, then the Kannada for friendship, then say \emph{bāndhavya} again
\end{itemize}

\begin{tcolorbox}[breakable,colback=teal!4,colframe=teal!35!black,title={Before you move on}]
What does \kn{ಬಾಂಧವ್ಯ} mean? (\textquotedblleft{}A bond; close ties\textquotedblright{}.) Say it once more.
\end{tcolorbox}
\section[\texorpdfstring{\kn{ಒಡನಾಟ} (oḍanāṭa)}{ಒಡನಾಟ (oḍanāṭa)}]{\kn{ಒಡನಾಟ} (oḍanāṭa) — companionship}

\label{lesson:KA-C528-odanata}



\begin{quote}
Before the new one: say the Kannada for friendship, then the Kannada for a bond.
\end{quote}

\subsection*{You'll want to know: \kn{ಒಡನಾಟ}}
\textbf{\kn{ಒಡನಾಟ}} — \emph{oḍanāṭa} — \textquotedblleft{}companionship\textquotedblright{}.

\begin{quote}
\kn{ಅಜ್ಜನ} \kn{ಒಡನಾಟದಲ್ಲಿ} \kn{ನಾನು} \kn{ತುಂಬಾ} \kn{ಕಲಿತೆ}.
\end{quote}

\begin{quote}
\emph{ajjana oḍanāṭadalli nānu tumbā kalite.}
\end{quote}

\begin{quote}
In grandfather's company I learned a great deal.
\end{quote}

Built on \kn{ಒಡನೆ}, together with.

\begin{cousinweb}
\textbf{\kn{ಒಡನಾಟ}} (\emph{oḍanāṭa}) is native Dravidian, built on Kannada \emph{oḍan}, \textquotedblleft{}union, together, with\textquotedblright{} (DEDR 945).
\end{cousinweb}

\begin{grammarlens}[title={What you've built}]
One more word for describing what you lived and felt.
\end{grammarlens}

\subsection*{Your turn}
\begin{itemize}
\raggedright
  \item \emph{Say it:} \emph{oḍanāṭa}
  \item \emph{Say it:} \emph{oḍanāṭa}, once more
  \item \emph{Say it:} \emph{ajjana oḍanāṭadalli nānu tumbā kalite.} — In grandfather's company I learned a great deal.
  \item \emph{Recall:} say the Kannada for friendship, then the Kannada for a bond, then say \emph{oḍanāṭa} again
\end{itemize}

\begin{tcolorbox}[breakable,colback=teal!4,colframe=teal!35!black,title={Before you move on}]
What does \kn{ಒಡನಾಟ} mean? (\textquotedblleft{}Companionship\textquotedblright{}.) Say it once more.
\end{tcolorbox}
\section[\texorpdfstring{\kn{ಆತ್ಮೀಯ} (ātmīya)}{ಆತ್ಮೀಯ (ātmīya)}]{\kn{ಆತ್ಮೀಯ} (ātmīya) — dear; close}

\label{lesson:KA-C528-atmiya}



\begin{quote}
Before the new one: say the Kannada for a bond, then the Kannada for companionship.
\end{quote}

\subsection*{You'll want to know: \kn{ಆತ್ಮೀಯ}}
\textbf{\kn{ಆತ್ಮೀಯ}} — \emph{ātmīya} — \textquotedblleft{}dear; close\textquotedblright{}.

\begin{quote}
\kn{ಅವನು} \kn{ನನ್ನ} \kn{ಆತ್ಮೀಯ} \kn{ಸ್ನೇಹಿತ}.
\end{quote}

\begin{quote}
\emph{avanu nanna ātmīya snēhita.}
\end{quote}

\begin{quote}
He is my close friend.
\end{quote}

Letters to friends begin \kn{ಆತ್ಮೀಯ} …, dear ….

\begin{cousinweb}
\textbf{\kn{ಆತ್ಮೀಯ}} (\emph{ātmīya}) is borrowed from Sanskrit \emph{ātmīya}, \textquotedblleft{}one's own\textquotedblright{} (Monier-Williams), built on \emph{ātman}, \textquotedblleft{}self\textquotedblright{}.
\end{cousinweb}

\begin{grammarlens}[title={What you've built}]
One more word for describing what you lived and felt.
\end{grammarlens}

\subsection*{Your turn}
\begin{itemize}
\raggedright
  \item \emph{Say it:} \emph{ātmīya}
  \item \emph{Say it:} \emph{ātmīya}, once more
  \item \emph{Say it:} \emph{avanu nanna ātmīya snēhita.} — He is my close friend.
  \item \emph{Recall:} say the Kannada for a bond, then the Kannada for companionship, then say \emph{ātmīya} again
\end{itemize}

\begin{tcolorbox}[breakable,colback=teal!4,colframe=teal!35!black,title={Before you move on}]
What does \kn{ಆತ್ಮೀಯ} mean? (\textquotedblleft{}Dear; close\textquotedblright{}.) Say it once more.
\end{tcolorbox}
\section[\texorpdfstring{\kn{ರಾಜಿ} \kn{ಮಾಡಿಕೊ} (rāji māḍiko)}{ರಾಜಿ ಮಾಡಿಕೊ (rāji māḍiko)}]{\kn{ರಾಜಿ} \kn{ಮಾಡಿಕೊ} (rāji māḍiko) — to make up; to reconcile}

\label{lesson:KA-C528-raji-madiko}



\begin{quote}
Before the new one: say the Kannada for companionship, then the Kannada for dear.
\end{quote}

\subsection*{You'll want to know: \kn{ರಾಜಿ} \kn{ಮಾಡಿಕೊ}}
\textbf{\kn{ರಾಜಿ} \kn{ಮಾಡಿಕೊ}} — \emph{rāji māḍiko} — \textquotedblleft{}to make up; to reconcile\textquotedblright{}.

\begin{quote}
\kn{ಜಗಳದ} \kn{ನಂತರ} \kn{ಇಬ್ಬರು} \kn{ಸ್ನೇಹಿತರು} \kn{ರಾಜಿ} \kn{ಮಾಡಿಕೊಂಡರು}.
\end{quote}

\begin{quote}
\emph{jagaḷada nantara ibbaru snēhitaru rāji māḍikoṇḍaru.}
\end{quote}

\begin{quote}
After the quarrel the two friends made up.
\end{quote}

Past (I …): \kn{ರಾಜಿ} \kn{ಮಾಡಿಕೊಂಡೆ}, \emph{rāji māḍikoṇḍe}.

With whom: -\kn{ಜೊತೆ}. \kn{ರಾಜಿ} is a settlement.

\begin{cousinweb}
\textbf{\kn{ರಾಜಿ} \kn{ಮಾಡಿಕೊ}} (\emph{rāji māḍiko}) joins \emph{rāji}, usually traced to Hindustani \emph{rāzī}, \textquotedblleft{}content, agreeing\textquotedblright{}, and \kn{ಮಾಡಿಕೊ}, \textquotedblleft{}to do for oneself\textquotedblright{}.
\end{cousinweb}

\begin{grammarlens}[title={What you've built}]
One more word for describing what you lived and felt.
\end{grammarlens}

\subsection*{Your turn}
\begin{itemize}
\raggedright
  \item \emph{Say it:} \emph{rāji māḍiko}
  \item \emph{Say it:} \emph{rāji māḍiko}, once more
  \item \emph{Say it:} \emph{jagaḷada nantara ibbaru snēhitaru rāji māḍikoṇḍaru.} — After the quarrel the two friends made up.
  \item \emph{Recall:} say the Kannada for companionship, then the Kannada for dear, then say \emph{rāji māḍiko} again
\end{itemize}

\begin{tcolorbox}[breakable,colback=teal!4,colframe=teal!35!black,title={Before you move on}]
What does \kn{ರಾಜಿ} \kn{ಮಾಡಿಕೊ} mean? (\textquotedblleft{}To make up; to reconcile\textquotedblright{}.) Say it once more.
\end{tcolorbox}
